\documentclass[11pt,letterpaper]{report}

% ---------------------------------------------------------------------
% Pacotes
% ---------------------------------------------------------------------
\usepackage[margin=1in]{geometry}
\usepackage[utf8]{inputenc}
\usepackage[T1]{fontenc}
\usepackage{xcolor}
\usepackage{listings}
\usepackage{longtable}
\usepackage{booktabs}
\usepackage{array}
\usepackage{enumitem}
\usepackage{hyperref}
\usepackage{fancyhdr}
\usepackage{titlesec}
\usepackage{tocloft}
\usepackage{xstring}

% ---------------------------------------------------------------------
% Cores
% ---------------------------------------------------------------------
\definecolor{cnimbusblue}{RGB}{25,55,110}
\definecolor{codebg}{RGB}{245,245,247}
\definecolor{codeframe}{RGB}{200,200,205}
\definecolor{commentgray}{RGB}{110,110,110}
\definecolor{keywordblue}{RGB}{20,90,160}

% ---------------------------------------------------------------------
% Estilo das seções
% ---------------------------------------------------------------------
\titleformat{\chapter}[display]
  {\normalfont\bfseries\color{cnimbusblue}}
  {\Large\chaptertitlename\ \thechapter}
  {10pt}{\Huge}
\titleformat{\section}
  {\normalfont\Large\bfseries\color{cnimbusblue}}{\thesection}{1em}{}
\titleformat{\subsection}
  {\normalfont\large\bfseries\color{cnimbusblue}}{\thesubsection}{1em}{}

\pagestyle{fancy}
\fancyhf{}
\fancyhead[L]{\small\textit{Manual do Usuário do cnimbus}}
\fancyhead[R]{\small\thepage}
\renewcommand{\headrulewidth}{0.4pt}

% ---------------------------------------------------------------------
% Blocos de código: dois estilos -- Nimbusfile (estilo Dockerfile) e shell
% ---------------------------------------------------------------------
\lstdefinelanguage{nimbusfile}{
  morekeywords={KERNEL,BUSYBOX,ARCH,VGA,HOSTNAME,DHCP,IP,DNS,NTP,FORMAT,
    USER,WORKDIR,LABEL,EXPOSE,ARG,ENV,VOLUME,AGENT,FIREWALL,FIREWALL6,
    FIREWALL_ON_ERROR,HEALTHCHECK,RESTART,STOPGRACE,TMPSIZE,PIECESKEY,
    COPY,ADD,ENTRYPOINT,CMD,SERVICE,HARDBOOT,WIFI,WIFIPSK,WIFICOUNTRY},
  sensitive=true,
  morecomment=[l]{\#},
  morestring=[b]",
}

\lstdefinestyle{nimbusstyle}{
  language=nimbusfile,
  backgroundcolor=\color{codebg},
  rulecolor=\color{codeframe},
  frame=single,
  framesep=4pt,
  basicstyle=\ttfamily\footnotesize,
  keywordstyle=\bfseries\color{keywordblue},
  commentstyle=\itshape\color{commentgray},
  stringstyle=\color{teal},
  breaklines=true,
  breakatwhitespace=false,
  columns=fullflexible,
  showstringspaces=false,
  tabsize=2,
  xleftmargin=1em,
  xrightmargin=1em,
}

\lstdefinestyle{shellstyle}{
  backgroundcolor=\color{black!92},
  basicstyle=\ttfamily\footnotesize\color{white},
  breaklines=true,
  breakatwhitespace=false,
  columns=fullflexible,
  frame=single,
  rulecolor=\color{black},
  framesep=4pt,
  showstringspaces=false,
  xleftmargin=1em,
  xrightmargin=1em,
}

\lstnewenvironment{nimbusfile}{\lstset{style=nimbusstyle}}{}
\lstnewenvironment{shell}{\lstset{style=shellstyle}}{}

% \flag/\dir/\cmdname renderizam flags de CLI e diretivas do Nimbusfile
% inline. \texttt{--foo} puro funde "--" em um en-dash sob T1/cm-super,
% então toda flag aparecia com um único traço -- corrigido substituindo
% "--" por "-{}-" (um grupo vazio quebra a fusão) antes de tipografar,
% uma vez só, na própria macro em vez de à mão em cada uma das muitas
% chamadas abaixo.
\newcommand{\dashfix}[1]{\StrSubstitute{#1}{--}{-{}-}}
\newcommand{\flag}[1]{\texttt{\textbf{\dashfix{#1}}}}
\newcommand{\dir}[1]{\texttt{\textbf{\dashfix{#1}}}}
\newcommand{\cmdname}[1]{\texttt{\textbf{\dashfix{#1}}}}

\hypersetup{
  colorlinks=true,
  linkcolor=cnimbusblue,
  urlcolor=cnimbusblue,
  citecolor=cnimbusblue,
  pdftitle={Manual do Usuário do cnimbus},
  pdfauthor={projeto cnimbus},
}

\setlength{\parskip}{0.5em}
\setlength{\parindent}{0pt}

% ---------------------------------------------------------------------
\begin{document}

% =======================================================================
% PÁGINA DE TÍTULO
% =======================================================================
\begin{titlepage}
  \centering
  \vspace*{2cm}
  {\Huge\bfseries\color{cnimbusblue} cnimbus}\\[0.5cm]
  {\Large Manual do Usuário}\\[0.3cm]
  {\large Construindo e Inicializando Imagens de VM Linux Mínimas e Distroless}\\[2cm]

  \rule{0.8\textwidth}{0.4pt}\\[0.4cm]
  {\large\itshape ``A mesma ideia de um Dockerfile: declare o que você quer, execute um comando.''}\\[0.4cm]
  \rule{0.8\textwidth}{0.4pt}\\[3cm]

  \begin{tabular}{rl}
    \textbf{Cobre:} & \texttt{prepare}, \texttt{build-disk}, \texttt{run}, \texttt{keygen} \\
    & Backends QEMU, VirtualBox, VMware, Hyper-V \\
    & A referência completa de diretivas do Nimbusfile \\
    & Assinatura de pieces, assinatura EFI-stub \& UKI \\
    & Boot em hardware físico (\texttt{HARDBOOT eth}/\texttt{wifi}/\texttt{eth+wifi}) \\
    & Exemplos práticos e compiláveis \\
  \end{tabular}

  \vfill
  {\small Esta versão do documento corresponde à linha de desenvolvimento v1.0 do \texttt{cnimbus}.\\
  Gerado como parte da própria documentação do projeto --
  veja \texttt{README.md} e \texttt{Tasks.md} na raiz do repositório
  para a fonte de verdade subjacente, continuamente validada, que
  este manual resume.\\[0.3cm]
  \textit{Termos técnicos, comandos, nomes de flags/diretivas e todo
  código permanecem em inglês, como no projeto original.}}
\end{titlepage}

\tableofcontents
\clearpage

% =======================================================================
\chapter{Introdução}
% =======================================================================

\section{O Que É o cnimbus}

\textbf{cnimbus} é uma ferramenta de linha de comando que constrói
imagens de máquina virtual Linux mínimas e com propósito específico a
partir de um único arquivo declarativo --- um \textbf{Nimbusfile} ---
e as inicializa sob um hypervisor real. O modelo mental é
deliberadamente o mesmo que o Docker popularizou: você declara o que
quer (uma versão de kernel, uma configuração de rede, um ou mais
processos para executar), e o \texttt{cnimbus} faz o resto: busca e
compila um kernel, monta um sistema de arquivos raiz mínimo baseado em
BusyBox, e monta uma imagem ISO ou disco raw inicializável.

Onde o cnimbus difere fortemente de um runtime de containers é no que
ele produz. Um container compartilha o kernel do host; uma imagem
cnimbus \textbf{é} seu próprio kernel mais um sistema de arquivos raiz
minúsculo (tipicamente abaixo de 20 MiB), somente leitura, residente em
RAM. Não há shell instalado por padrão, nenhum gerenciador de pacotes,
e nenhum armazenamento persistente gravável a menos que você declare
explicitamente um \dir{VOLUME}. Cada imagem é construída para exatamente
um propósito: executar o processo (ou o punhado de processos) que você
declarou, e nada mais.

\section{Filosofia de Design}

Quatro decisões moldam todo o resto deste manual:

\begin{itemize}
  \item \textbf{Distroless.} Não há um userland de propósito geral.
    O BusyBox fornece apenas o pequeno conjunto de applets que a
    sequência de boot e os \dir{SERVICE}s que você declarou realmente
    precisam (\texttt{busybox -{}-install} conecta exatamente esses,
    não os cerca de 400 applets que uma build padrão do BusyBox traz).
  \item \textbf{Sem publicação, sem registro hospedado.} O cnimbus não
    hospeda nem distribui ``pieces'' pré-construídas (binários de
    kernel/BusyBox/iptables). Você constrói as suas próprias, uma vez,
    e as armazena em cache localmente ou em infraestrutura que você
    controla. Não existe um equivalente do Docker Hub para o cnimbus,
    por decisão explícita de produto.
  \item \textbf{Somente auto-construído.} Todo artefato que termina na
    imagem é ou compilado pelo \texttt{cnimbus prepare} a partir de
    fonte upstream (kernel.org, BusyBox, iptables, wpa\_supplicant) com
    sua proveniência verificada (assinaturas PGP onde o upstream
    publica, hash fixado em todo o resto), ou explicitamente copiado
    por você via \dir{COPY}. Nada é buscado de terceiros não
    verificados em tempo de boot ou de build.
  \item \textbf{Comportamento real, validado por boot.} Todo mecanismo
    descrito neste manual --- cada diretiva, cada backend, cada flag
    --- foi exercitado contra uma máquina virtual real, de verdade
    inicializada, durante o desenvolvimento deste projeto, não apenas
    testado por unidade. Onde uma limitação real foi encontrada (e
    várias foram), ela é documentada como tal, em vez de disfarçada.
\end{itemize}

\section{Como Este Manual Está Organizado}

\begin{itemize}
  \item O Capítulo~\ref{ch:install} coloca o \texttt{cnimbus} compilado
    e suas primeiras pieces prontas.
  \item O Capítulo~\ref{ch:quickstart} percorre o caminho mais curto
    possível de um diretório vazio até uma VM inicializada.
  \item O Capítulo~\ref{ch:nimbusfile} é a referência completa de
    diretivas do Nimbusfile --- o coração deste manual.
  \item O Capítulo~\ref{ch:cli} documenta todo subcomando e flag.
  \item O Capítulo~\ref{ch:backends} cobre os quatro backends de
    hypervisor suportados e como alcançar um serviço rodando dentro do
    guest.
  \item O Capítulo~\ref{ch:security} cobre a assinatura de pieces e a
    assinatura de kernel EFI-stub / Secure Boot / UKI.
  \item O Capítulo~\ref{ch:hardboot} cobre o perfil de boot opt-in para
    hardware físico (bare-metal).
  \item O Capítulo~\ref{ch:examples} percorre cada exemplo enviado no
    diretório \texttt{examples/} do repositório, do início ao fim.
  \item O Capítulo~\ref{ch:troubleshooting} reúne limitações conhecidas
    e suas soluções alternativas.
  \item Os apêndices trazem uma referência rápida de uma página das
    diretivas e as métricas medidas de requisitos não-funcionais do
    projeto.
\end{itemize}

% =======================================================================
\chapter{Instalação}
\label{ch:install}
% =======================================================================

\section{Requisitos}

\begin{itemize}
  \item \textbf{Docker} (Desktop no Windows/macOS, Engine no Linux) ou
    \textbf{Podman} --- este é o único requisito real. É onde a
    compilação real de kernel/BusyBox/iptables acontece
    (\texttt{cnimbus prepare}); \texttt{cnimbus build-disk} e
    \texttt{cnimbus run} não precisam de Docker/Podman nenhum, uma vez
    que as pieces já existam. Go \textit{não} é um requisito: o próprio
    binário \texttt{cnimbus} pode ser compilado inteiramente dentro de
    um container Docker/Podman, sem nenhum Go instalado no host --- veja
    \texttt{BUILD.md} na raiz do repositório.
  \item \textbf{Opcional} --- só necessário para de fato inicializar
    uma imagem com \texttt{cnimbus run} (o objetivo central do projeto
    é montar a imagem; inicializá-la é um passo posterior e dispensável):
    um entre \textbf{QEMU}, \textbf{VirtualBox}, \textbf{VMware
    Workstation/Player}, ou \textbf{Hyper-V}.
\end{itemize}

\section{Compilando a Partir do Código-Fonte}

Com Go instalado no host:

\begin{shell}
git clone <este-repositorio>
cd cnimbus
go build -o cnimbus ./cmd/cnimbus
\end{shell}

Sem Go instalado no host --- apenas Docker (ou Podman), o único
requisito real do projeto:

\begin{shell}
docker run --rm -v "$(pwd)":/src -w /src \
  -e CGO_ENABLED=0 -e GOOS=linux -e GOARCH=amd64 \
  golang:1.26.5 go build -o cnimbus ./cmd/cnimbus
\end{shell}

Ambos os caminhos produzem um único binário \texttt{cnimbus} (ou
\texttt{cnimbus.exe} no Windows) sem dependências. Veja
\texttt{BUILD.md} na raiz do repositório para a matriz completa de
compilação cruzada (Windows, Linux e macOS, em amd64, arm64 e ---
experimentalmente --- riscv64 só para Linux), incluindo os comandos
equivalentes para cada plataforma via Docker/Podman.

\section{Compilando Suas Primeiras Pieces}

\texttt{prepare} é o único subcomando que fala com o Docker. Ele baixa
e compila um kernel (verificando sua assinatura PGP contra as chaves
publicadas pelo próprio kernel.org, buscadas em tempo real via Web Key
Directory), constrói um BusyBox correspondente, e constrói um
\texttt{iptables} estático --- juntos chamados de \textit{pieces}.

\begin{shell}
./cnimbus prepare --out ./pieces
\end{shell}

A primeira execução leva alguns minutos (uma compilação real de
kernel). Toda execução posterior é instantânea se nada mudou, já que o
Thunder (o pequeno programa em Go que conduz a build real dentro do
container Docker) verifica um lockfile sha256 antes de decidir pular o
trabalho. O resultado fica namespaced por arquitetura em
\texttt{./pieces/amd64/} (ou \texttt{./pieces/arm64/}):

\begin{shell}
pieces/amd64/
  vmlinuz
  busybox
  busybox-manifest.tsv
  iptables
  pieces.json
  pieces.sha256
\end{shell}

Tudo a partir daqui --- todo exemplo neste manual --- assume que esse
diretório existe.

% =======================================================================
\chapter{Começo Rápido}
\label{ch:quickstart}
% =======================================================================

Este capítulo constrói a menor imagem funcional possível: um servidor
HTTP estático em Go, alcançável a partir do seu host, inicializado sob
o QEMU.

\section{Passo 1: Um Servidor Mínimo}

\begin{shell}
mkdir my-first-vm && cd my-first-vm
cat > server.go << 'EOF'
package main

import ("fmt"; "log"; "net/http")

func main() {
    http.HandleFunc("/", func(w http.ResponseWriter, r *http.Request) {
        fmt.Fprintln(w, "hello from cnimbus")
    })
    log.Println("listening on :8080")
    log.Fatal(http.ListenAndServe(":8080", nil))
}
EOF
GOOS=linux GOARCH=amd64 CGO_ENABLED=0 go build -o helloserver server.go
\end{shell}

\texttt{CGO\_ENABLED=0} importa: a imagem não tem biblioteca C
instalada, então um binário dinamicamente ligado falharia ao iniciar.

\section{Passo 2: O Nimbusfile}

Um Nimbusfile mínimo usando os valores padrão (\texttt{latest-stable},
\texttt{latest}) já é suficiente para o exemplo acima. Mas em builds
reais e reproduzíveis --- especialmente em CI, ou quando você precisa
que o mesmo Nimbusfile produza bit a bit a mesma imagem meses depois
--- é melhor fixar versões exatas em vez de confiar em ``a mais recente
de hoje''. Os dois exemplos abaixo mostram ambos os estilos.

\textbf{Nimbusfile mínimo (versões flutuantes):}

\begin{nimbusfile}
KERNEL latest-stable
BUSYBOX latest
ARCH amd64
HOSTNAME my-first-vm
DHCP true

COPY ./helloserver /usr/bin/helloserver
ENTRYPOINT /usr/bin/helloserver
CMD :8080

FORMAT iso
\end{nimbusfile}

\textbf{Nimbusfile com versões reais fixadas (build reproduzível):}

\begin{nimbusfile}
KERNEL 6.9.4
BUSYBOX 1.36.1
ARCH amd64
HOSTNAME my-first-vm-pinned
DHCP true

COPY ./helloserver /usr/bin/helloserver
ENTRYPOINT /usr/bin/helloserver
CMD :8080

FORMAT iso
\end{nimbusfile}

Com \dir{KERNEL}/\dir{BUSYBOX} fixados em números de versão exatos (em
vez de \texttt{latest-stable}/\texttt{latest}), o mesmo Nimbusfile
compila exatamente o mesmo par kernel+BusyBox hoje e em qualquer data
futura --- o único jeito de garantir que uma imagem seja realmente
reproduzível, já que ``a versão mais recente do kernel.org'' muda com
o tempo por definição. Antes de testar que este segundo Nimbusfile
realmente funciona, compile as pieces fixadas e depois monte a imagem:

\begin{shell}
../cnimbus prepare --kernel 6.9.4 --busybox 1.36.1 --out ../pieces-pinned
../cnimbus build-disk -f Nimbusfile.pinned --pieces ../pieces-pinned
\end{shell}

\section{Passo 3: Compilar}

\begin{shell}
../cnimbus build-disk --pieces ../pieces
\end{shell}

Isso produz \texttt{my-first-vm.iso} no diretório atual --- tipicamente
com menos de 20 MiB.

\section{Passo 4: Inicializar e Alcançar}

\begin{shell}
../cnimbus run my-first-vm.iso
\end{shell}

Por padrão isso inicializa sob o QEMU com um console serial e
redireciona a porta 8080 do host para a porta 8080 do guest. De outro
terminal:

\begin{shell}
curl http://127.0.0.1:8080/
# hello from cnimbus
\end{shell}

Esse é o ciclo completo: declarar, compilar, inicializar, alcançar.
Todo o restante deste manual trata das muitas coisas que você pode
adicionar ao Nimbusfile, ou flags que você pode passar para mudar como
a imagem é construída ou inicializada.

% =======================================================================
\chapter{O Nimbusfile}
\label{ch:nimbusfile}
% =======================================================================

\section{Conceitos}

Um Nimbusfile é um arquivo de texto puro, uma diretiva por linha
(linhas longas podem ser continuadas com uma barra invertida no final,
como em um Dockerfile). Comentários começam com \texttt{\#}.
\texttt{cnimbus init} escreve um Nimbusfile inicial totalmente
comentado no diretório atual --- útil como referência viva ao lado
deste capítulo.

Existem duas categorias de diretiva:

\begin{itemize}
  \item \textbf{Lidas apenas pelo \texttt{prepare}} (\dir{KERNEL},
    \dir{BUSYBOX}, \dir{VGA}, \dir{HARDBOOT}): afetam qual kernel é
    compilado, então precisam ser conhecidas \textit{antes} das pieces
    existirem, não apenas antes da imagem ser montada.
  \item \textbf{Lidas apenas pelo \texttt{build-disk}} (todo o resto):
    afetam como a imagem é montada a partir de pieces já existentes, e
    podem mudar entre builds sem recompilar nada.
\end{itemize}

\dir{ARCH} é lida por \textit{ambos} os comandos, já que determina
tanto o que o \texttt{prepare} compila quanto quais pieces
namespaced-por-arquitetura o \texttt{build-disk} deve buscar.

Toda diretiva que também tem uma flag de linha de comando com o mesmo
nome segue uma regra consistente: \textbf{a flag sempre vence} sobre o
Nimbusfile. Isso permite que um único Nimbusfile sirva a múltiplas
builds (uma matriz de CI construindo tanto \texttt{amd64} quanto
\texttt{arm64} a partir do mesmo arquivo, por exemplo) sem editá-lo.

\section{Referência Completa de Diretivas}

\subsection{Identidade da Imagem e Alvo de Boot}

\begin{longtable}{p{4.6cm}p{8.9cm}}
\toprule
\textbf{Diretiva} & \textbf{Significado} \\
\midrule
\endhead
\dir{KERNEL <version>} & \texttt{latest-stable}, \texttt{latest-longterm}, ou uma versão explícita (ex. \texttt{6.9.4}). Lida apenas pelo \texttt{prepare}. \\
\dir{BUSYBOX <version>} & Uma versão explícita, ou \texttt{latest} para o padrão do próprio cnimbus. Lida apenas pelo \texttt{prepare}. \\
\dir{ARCH <amd64|arm64>} & Arquitetura de destino (padrão \texttt{amd64}). Lida tanto pelo \texttt{prepare} quanto pelo \texttt{build-disk}. \\
\dir{VGA <true|false>} & Habilita um console VGA/framebuffer real (padrão \texttt{false} --- o console serial está sempre disponível independentemente). Também ativa um banner no console: assim que a rede sobe, todo endereço IPv4 (e IPv6, se houver) de cada interface é impresso na tela --- a única forma de descobrir o endereço da imagem sem um shell para entrar. Lida apenas pelo \texttt{prepare}. \\
\dir{HARDBOOT <none|eth|wifi|eth+wifi>} & Perfil de boot em hardware físico (padrão \texttt{none}). Muda quais drivers de kernel são compilados --- veja o Capítulo~\ref{ch:hardboot}. Lida apenas pelo \texttt{prepare}. \\
\dir{FORMAT <iso|raw|vhd>} & \texttt{iso}: uma imagem de disco óptico inicializável para uso em VM. \texttt{raw}: um disco GPT (ESP UEFI + partição raiz SquashFS), o formato para templates de disco em nuvem \textit{e} boot USB/bare-metal. \texttt{vhd}: o mesmo layout raw envolto em um footer VHD Fixed real, pronto para o Hyper-V sem precisar de um passo \texttt{run} separado. \\
\bottomrule
\end{longtable}

\subsection{Rede}

\begin{longtable}{p{4.2cm}p{9.3cm}}
\toprule
\textbf{Diretiva} & \textbf{Significado} \\
\midrule
\endhead
\dir{HOSTNAME <name>} & Hostname da imagem. \\
\dir{DHCP <true|false>} & Sobe a \texttt{eth0} via DHCP no boot. \\
\dir{IP <addr> <mask> <gw>} & IP estático em vez de DHCP; vence o DHCP se ambos estiverem definidos. \\
\dir{DNS <addr>} & Um nameserver explícito, sobrescrevendo o que o DHCP forneceu. Repetível. \\
\dir{NTP <server|false>} & Sincroniza o relógio no boot (padrão \texttt{pool.ntp.org}). Repetível; \texttt{false} desativa. Precisa de rede configurada. \\
\dir{WIFI <ssid>} & Rede WiFi para associar. Requer \dir{HARDBOOT wifi} ou \dir{HARDBOOT eth+wifi}. Veja o Capítulo~\ref{ch:hardboot}. \\
\dir{WIFIPSK <passphrase|hex-key>} & Credencial WPA2-PSK para \dir{WIFI}. Requer \dir{HARDBOOT wifi} ou \dir{HARDBOOT eth+wifi}. \\
\dir{WIFICOUNTRY <code>} & Domínio regulatório de duas letras (ex. \texttt{US}, \texttt{BR}) para o rádio WiFi. Requer \dir{HARDBOOT wifi} ou \dir{HARDBOOT eth+wifi}. \\
\bottomrule
\end{longtable}

\subsection{Processo e Sistema de Arquivos}

\begin{longtable}{p{4.2cm}p{9.3cm}}
\toprule
\textbf{Diretiva} & \textbf{Significado} \\
\midrule
\endhead
\dir{USER <name>} & Executa todo processo com esta conta sem privilégios (uid/gid 1000) em vez de root. \\
\dir{WORKDIR <path>} & Diretório de trabalho para todo processo (padrão \texttt{/}). \\
\dir{ENV <KEY>=<VALUE>} & Variável de ambiente exportada para todo processo. Repetível. \\
\dir{LABEL <KEY>=<VALUE>} & Metadado livre escrito em \texttt{/etc/cnimbus-release}. Sem efeito no boot. Repetível. \\
\dir{EXPOSE <port>[/tcp|/udp]} & Documenta uma porta de escuta. Apenas informativo. Repetível. \\
\dir{ARG <NAME>[=<default>]} & Variável de tempo de build, usável como \texttt{\$\{NAME\}} em outros lugares. Sobrescreva com \flag{--build-arg NAME=value}. \\
\dir{COPY [--chmod=<mode>] <src> <dest>} & Arquivo/diretório/glob local para dentro da imagem. Deve corresponder ao \dir{ARCH}. \\
\dir{ADD [--chmod=<mode>] <src> <dest>} & Como \dir{COPY}, mas \texttt{src} pode ser uma URL, e arquivos compactados são auto-extraídos. \\
\dir{ENTRYPOINT <cmd...>} & O processo principal, reiniciado automaticamente. Forma shell ou exec. \\
\dir{CMD <args...>} & Argumentos padrão anexados após \dir{ENTRYPOINT}. \\
\dir{SERVICE <name> <cmd...>} & Um processo adicional, supervisionado e reiniciado automaticamente. Repetível. \\
\bottomrule
\end{longtable}

\subsection{Armazenamento}

\begin{longtable}{p{4.2cm}p{9.3cm}}
\toprule
\textbf{Diretiva} & \textbf{Significado} \\
\midrule
\endhead
\dir{VOLUME <device> <mount> [fstype] [required]} & Monta um disco pré-existente e pré-formatado no boot. Nunca o formata. \texttt{fstype} é \texttt{vfat} (padrão) ou \texttt{ext4}. Sem \texttt{required}: uma montagem que falha é silenciosamente ignorada. Com \texttt{required}: uma montagem que falha interrompe o boot com uma mensagem FATAL clara. Repetível. \\
\dir{TMPSIZE <size>} & Sobrescreve o \texttt{size=} do tmpfs para os diretórios executáveis em RAM (\texttt{bin}, \texttt{sbin}, \texttt{usr/bin}, \texttt{usr/sbin}); padrão \texttt{32m}. \\
\bottomrule
\end{longtable}

\subsection{Configuração ao Vivo --- A Diretiva AGENT}

A família \dir{AGENT} permite que uma VM em execução capte mudanças de
configuração sem um rebuild ou reboot, escrevendo o valor buscado em
\texttt{/var/run/cnimbus-kv.json} em um intervalo que seu
\dir{ENTRYPOINT} lê.

\begin{longtable}{p{5.4cm}p{8.1cm}}
\toprule
\textbf{Forma} & \textbf{Canal} \\
\midrule
\endhead
\dir{AGENT <url> [interval]} & HTTP(S) puro, qualquer servidor. Funciona em qualquer hypervisor com rede no guest. Intervalo padrão \texttt{5}s. \\
\dir{AGENT header <name>: <value>} & Cabeçalho extra para a forma HTTP acima (endpoints de metadados em nuvem frequentemente exigem um). Repetível. \\
\dir{AGENT vboxguest <property> [interval]} & O canal real de Guest Properties do VirtualBox. Não precisa de rede no guest. Exclusivo do VirtualBox. \\
\dir{AGENT virtio-serial <device> [interval]} & Um dispositivo de caractere \texttt{virtio-serial} do QEMU/Proxmox. Não precisa de \texttt{qemu-guest-agent}. \\
\dir{AGENT aws-imds <path> [interval]} & Serviço de metadados IMDSv2 do AWS EC2 (trata o handshake de PUT-token-depois-GET). \\
\dir{AGENT ibm-imds <path> [interval]} & Serviço de metadados do IBM Cloud VPC. \\
\dir{AGENT vmware <key> [interval]} & \texttt{guestinfo.<key>} do VMware via o protocolo de I/O backdoor de baixa banda. Não precisa das VMware Tools. Só \texttt{linux/amd64}. \\
\bottomrule
\end{longtable}

\subsection{Firewall e Ciclo de Vida}

\begin{longtable}{p{4.2cm}p{9.3cm}}
\toprule
\textbf{Diretiva} & \textbf{Significado} \\
\midrule
\endhead
\dir{FIREWALL <rule>} & Uma linha de regra \texttt{iptables} (IPv4), aplicada no boot via um \texttt{iptables} estático construído pelo \texttt{prepare}. Repetível. \\
\dir{FIREWALL6 <rule>} & O equivalente IPv6 de \dir{FIREWALL}: mesma sintaxe de regra, aplicada via o mesmo binário estático (que também despacha como \texttt{ip6tables}), num conjunto de regras totalmente independente. Declarar um não afeta o outro. Repetível. Injeta automaticamente as mensagens ICMPv6 mínimas da RFC~4890 (Neighbor Discovery e erros de MTU) antes de qualquer regra sua --- sem isso, uma política \texttt{-P INPUT DROP} bloquearia a própria resolução de endereço IPv6 (o equivalente do ARP), algo que \dir{FIREWALL} nunca precisa porque o ARP não passa pelo \texttt{iptables}. \\
\dir{FIREWALL\_ON\_ERROR <open|closed>} & Comportamento de fallback se uma regra \dir{FIREWALL}/\dir{FIREWALL6} falhar ao ser aplicada: \texttt{open} (padrão, aceita tudo) ou \texttt{closed} (bloqueia tudo exceto loopback/estabelecidas). Vale para os dois conjuntos de regras. \\
\dir{HEALTHCHECK [--interval=<n>] [--retries=<n>] <cmd...>} & Executa \texttt{<cmd...>} periodicamente contra o \dir{ENTRYPOINT}; mata/reinicia após falhas consecutivas. Padrões: intervalo \texttt{30}s, tentativas \texttt{3}. \\
\dir{RESTART <target> <always|on-failure|no>} & Política de reinício para \texttt{entrypoint} ou um \dir{SERVICE} nomeado. \\
\dir{STOPGRACE <seconds>} & Quanto tempo um shutdown espera o \dir{ENTRYPOINT} sair por conta própria após \texttt{SIGTERM} antes do \texttt{SIGKILL}. Padrão \texttt{10}. \\
\bottomrule
\end{longtable}

\subsection{Assinatura}

\begin{longtable}{p{4.2cm}p{9.3cm}}
\toprule
\textbf{Diretiva} & \textbf{Significado} \\
\midrule
\endhead
\dir{PIECESKEY <hex-pubkey>} & Fixa a chave pública Ed25519 com a qual o \texttt{pieces.sha256} deve estar assinado. Veja o Capítulo~\ref{ch:security}. \\
\bottomrule
\end{longtable}

\section{Regras de Precedência Que Vale Memorizar}

\begin{enumerate}
  \item Uma flag de linha de comando sempre sobrescreve a diretiva de
    mesmo nome no Nimbusfile.
  \item \dir{IP} vence \dir{DHCP} quando ambos estão presentes.
  \item Um \dir{ENV} posterior com a mesma chave sobrescreve um
    anterior.
  \item \texttt{\$\$} é o escape para um \texttt{\$} literal em
    qualquer argumento de diretiva (convenção do próprio Docker) ---
    necessário para que \dir{ENTRYPOINT}/\dir{CMD} passem um
    \texttt{\$VAR} literal adiante, para o shell expandir em
    \textit{tempo de execução}, a partir de um valor \dir{ENV}, em vez
    de em tempo de \textit{parse} do Nimbusfile.
\end{enumerate}

% =======================================================================
\chapter{Referência de Linha de Comando}
\label{ch:cli}
% =======================================================================

\section{cnimbus prepare}

Compila as pieces (kernel, BusyBox, iptables estático) dentro do
Docker.

\begin{longtable}{p{5.3cm}p{8.2cm}}
\toprule
\textbf{Flag} & \textbf{Significado} \\
\midrule
\endhead
\flag{-f <path>} & Nimbusfile de onde ler os padrões de \dir{KERNEL}/\dir{BUSYBOX}/\dir{ARCH} (padrão \texttt{Nimbusfile}). \\
\flag{--kernel <version>} & Sobrescreve o \dir{KERNEL} do Nimbusfile. \\
\flag{--busybox <version>} & Sobrescreve o \dir{BUSYBOX} do Nimbusfile. \\
\flag{--arch <amd64|arm64>} & Sobrescreve o \dir{ARCH} do Nimbusfile. \\
\flag{--out <dir>} & Diretório onde escrever as pieces (padrão \texttt{./pieces}, namespaced por arquitetura). \\
\flag{--vga[=false]} & Sobrescreve o \dir{VGA} do Nimbusfile. \\
\flag{--hardboot <none|eth|wifi|eth+wifi>} & Sobrescreve o \dir{HARDBOOT} do Nimbusfile. Veja o Capítulo~\ref{ch:hardboot}. \\
\flag{--insecure-skip-kernel-verify} & Pula a verificação PGP do tarball do kernel baixado. Não recomendado. \\
\flag{--pieces-sign-key <path>} & Assina o \texttt{pieces.sha256} recém-escrito com esta chave Ed25519. \\
\flag{--jobs <n>} & Paralelismo \texttt{make -j} para as builds de kernel/BusyBox/iptables (0 = auto-detecta a partir da memória disponível). \\
\bottomrule
\end{longtable}

\section{cnimbus build-disk}

Monta uma imagem inicializável a partir de pieces já construídas. Nunca
toca no Docker.

\begin{longtable}{p{5.3cm}p{8.2cm}}
\toprule
\textbf{Flag} & \textbf{Significado} \\
\midrule
\endhead
\flag{-f <path>} & Nimbusfile a ler (padrão \texttt{Nimbusfile}). \\
\flag{-o <path>} & Caminho da imagem de saída (padrão \texttt{<hostname>.iso}, \texttt{.img} ou \texttt{.vhd}). \\
\flag{--pieces <src>} & Fonte de pieces pré-construídas: um diretório local ou um prefixo de URL \texttt{http(s)://}. \\
\flag{--arch <amd64|arm64>} & Sobrescreve o \dir{ARCH} do Nimbusfile. \\
\flag{--pieces-verify-key <hex>} & Uma chave pública Ed25519 contra a qual o \texttt{pieces.sha256.sig} deve verificar. \\
\flag{--pieces-cache-dir <dir>} & Cache local para uma fonte \flag{--pieces} remota. \\
\flag{--no-pieces-cache} & Desativa completamente o cache de pieces. \\
\flag{--pieces-insecure-http} & Permite uma fonte \flag{--pieces} \texttt{http://} pura (recusada por padrão). \\
\flag{--pieces-allow-unverified} & Permite uma fonte de pieces sem assinatura, sem manifesto de integridade nenhum. \\
\flag{--no-lockfile} & Não escreve o \texttt{<output>.lock}. \\
\flag{--tmpdir <dir>} & Diretório para os arquivos temporários de trabalho da build. \\
\flag{--secureboot} & Assina o kernel EFI-stub enviado com uma implementação Authenticode em Go puro (sem Docker). Veja o Capítulo~\ref{ch:security}. \\
\flag{--uki} & Monta e assina uma Unified Kernel Image (implica \flag{--secureboot}). \\
\flag{--secureboot-key <path>} & Traz sua própria chave privada RSA para assinatura (deve vir junto com \flag{--secureboot-cert}). \\
\flag{--secureboot-cert <path>} & Traz seu próprio certificado X.509 para assinatura. \\
\flag{--secureboot-dir <dir>} & Onde auto-gerar/reutilizar um par de chaves de assinatura quando nenhuma das opções acima é dada (padrão \texttt{./secureboot}). \\
\bottomrule
\end{longtable}

\section{cnimbus run}

Inicializa uma imagem já construída sob um hypervisor real.

\begin{longtable}{p{5.3cm}p{8.2cm}}
\toprule
\textbf{Flag} & \textbf{Significado} \\
\midrule
\endhead
\flag{--backend <qemu|vbox|vmware|hyperv>} & Hypervisor a usar (padrão \texttt{qemu}). \\
\flag{--arch <amd64|arm64>} & Arquitetura da imagem (padrão \texttt{amd64}). \\
\flag{--uefi} & Inicializa via UEFI/OVMF em vez de BIOS (só amd64 --- arm64 é sempre UEFI). \\
\flag{--ovmf-code <path>} / \flag{--ovmf-vars <path>} & Caminhos explícitos do firmware OVMF (auto-detectados caso contrário). \\
\flag{--mem <MB>} & RAM do guest (padrão \texttt{512}). \\
\flag{--smp <n>} & Número de vCPUs do guest (só \flag{--backend qemu}). \\
\flag{--accel <auto|kvm|hvf|whpx|tcg>} & Backend de aceleração do QEMU (só \flag{--backend qemu}). \\
\flag{--hostfwd <host:guest>} & Redirecionamento de porta TCP (padrão \texttt{8080:8080}). \\
\flag{--hostfwd-bind <addr>} & Endereço do host ao qual o redirecionamento se vincula (padrão \texttt{127.0.0.1}). \\
\flag{--backend vbox|vmware|hyperv} & Veja o Capítulo~\ref{ch:backends} para notas específicas de cada backend. \\
\flag{--vm-name <name>} & Nome da VM a criar (removida após o desligamento a menos que \flag{--vm-keep}). \\
\flag{--vm-keep} & Não remove a VM após ela desligar. \\
\bottomrule
\end{longtable}

\section{cnimbus keygen}

Gera as chaves de assinatura usadas em outros pontos deste manual.

\begin{shell}
cnimbus keygen --out my-signing-key.hex          # assinatura de pieces (Ed25519)
cnimbus keygen --secureboot --out-dir ./secureboot # Secure Boot (RSA + X.509)
\end{shell}

\section{Outros Subcomandos}

\begin{itemize}
  \item \cmdname{cnimbus init} --- escreve um Nimbusfile inicial
    totalmente comentado no diretório atual.
  \item \cmdname{cnimbus validate} --- verifica um Nimbusfile (e a
    arquitetura de qualquer binário copiado via \dir{COPY}) sem
    construir nada.
  \item \cmdname{cnimbus kv-serve} --- um pequeno servidor HTTP local
    para testar o mecanismo \dir{AGENT <url>}, com TLS opcional e
    autenticação por bearer token.
  \item \cmdname{cnimbus help} --- lista todo subcomando e os códigos
    de saída documentados do projeto.
\end{itemize}

% =======================================================================
\chapter{Backends de Hypervisor}
\label{ch:backends}
% =======================================================================

\section{Visão Geral}

\texttt{cnimbus run} unifica quatro backends sob uma única flag,
\flag{--backend}. A Tabela~\ref{tab:backends} resume o que cada um
precisa e suporta.

\begin{table}[h]
\centering
\begin{tabular}{lllll}
\toprule
\textbf{Backend} & \textbf{Precisa} & \textbf{UEFI} & \textbf{FORMAT raw} & \textbf{Redir. porta} \\
\midrule
\texttt{qemu} & QEMU no PATH & Sim (OVMF) & Sim & \texttt{-netdev user} \\
\texttt{vbox} & VBoxManage & Sim & Sim & Redir. NAT \\
\texttt{vmware} & vmrun/VMware & Sim & Sim (envolto em VMDK) & Não automatizado \\
\texttt{hyperv} & Windows, admin & Sim (Ger.\,2) & Só sim & Switch NAT \\
\bottomrule
\end{tabular}
\caption{Resumo de capacidades dos backends.}
\label{tab:backends}
\end{table}

\section{QEMU}

O padrão. Nenhuma elevação necessária em qualquer plataforma.

\begin{shell}
cnimbus run my-image.iso
cnimbus run --uefi --accel kvm my-image.iso   # Linux com KVM
\end{shell}

\section{VirtualBox}

\begin{shell}
cnimbus run --backend vbox my-image.iso
cnimbus run --backend vbox --uefi my-image.iso
\end{shell}

Nenhuma elevação é necessária em nenhuma plataforma testada por este
projeto --- o \texttt{VBoxManage} executa diretamente como o usuário
que o invoca.

\section{VMware}

\begin{shell}
cnimbus run --backend vmware my-image.iso
\end{shell}

O VMware Player (a versão gratuita) não pode usar vTPM --- o Secure
Boot funciona, mas a emulação de TPM exige o Workstation Pro ou o
vSphere (criptografia da VM é um pré-requisito obrigatório que o
Player não expõe). O redirecionamento de porta não é automatizado
para este backend; use a própria configuração de rede do guest (um
adaptador em bridge/NAT que você configura diretamente no VMware) para
alcançar um serviço rodando dentro dele.

\section{Hyper-V}

\begin{shell}
cnimbus run --backend hyperv --uefi my-image.img
\end{shell}

\textbf{Exige \flag{FORMAT raw}} --- o Hyper-V não tem um caminho de
drive óptico virtual na implementação deste projeto, apenas um disco
rígido virtual (um VHD Fixed, envolvendo a imagem raw de forma
transparente). Exige uma sessão elevada (Administrador): execute a
partir de uma janela do PowerShell aberta via ``Executar como
administrador'', ou aprove o prompt do UAC se o próprio
\texttt{cnimbus} disparar um. Um guest alcançável via \flag{--hostfwd}
deve declarar um \dir{IP} estático na faixa \texttt{192.168.200.0/24}
--- o switch NAT interno que este backend cria não tem servidor DHCP
próprio.

% =======================================================================
\chapter{Segurança: Assinatura e Secure Boot}
\label{ch:security}
% =======================================================================

\section{Assinatura de Pieces (Integridade e Autenticidade)}

O \texttt{pieces.sha256} (escrito pelo \texttt{prepare}) é um manifesto
de hashes: o \texttt{build-disk} verifica cada arquivo buscado contra
ele, então um download corrompido ou uma substituição em trânsito é
detectada. Isso é \textit{integridade} --- não prova \textit{quem}
publicou o manifesto.

\begin{shell}
cnimbus keygen --out my-signing-key.hex
cnimbus prepare --pieces-sign-key my-signing-key.hex
cnimbus build-disk --pieces-verify-key <a-chave-publica-impressa>
\end{shell}

Com \flag{--pieces-verify-key} (ou uma linha \dir{PIECESKEY} no
Nimbusfile) definida, o \texttt{build-disk} se recusa a construir a
menos que a fonte de pieces tenha publicado um \texttt{pieces.sha256.sig}
que verifique contra exatamente essa chave.

\section{Assinatura de Kernel EFI-Stub e UKI}

Assinar as pieces autentica as \textit{entradas da build}. Isso não
diz nada sobre o que o \textit{firmware} confia em tempo de boot.
\flag{--secureboot} e \flag{--uki} fecham essa lacuna: elas assinam a
imagem EFI de boot real com uma implementação Authenticode em Go puro
(sem Docker, sem ferramenta externa nenhuma --- veja
\texttt{internal/secureboot}), então um hypervisor com Secure Boot
habilitado vai se recusar a executá-la a menos que seu certificado
esteja registrado no seu banco de dados UEFI.

\begin{shell}
cnimbus build-disk --secureboot          # assina o kernel diretamente
cnimbus build-disk --uki                 # monta + assina uma UKI completa
\end{shell}

Sem flags explícitas de chave/certificado, uma identidade de assinatura
(RSA-3072 + X.509 autoassinado) é gerada uma vez sob
\flag{--secureboot-dir} (padrão \texttt{./secureboot}) e reutilizada em
toda build posterior --- nunca regenerada silenciosamente. Veja o
passo a passo de \texttt{examples/secureboot-uki/} no
Capítulo~\ref{ch:examples} para os comandos exatos, testados de ponta a
ponta, para registrar esse certificado na variável \texttt{db} UEFI de
uma VM VirtualBox real e inicializar com o Secure Boot habilitado.

\textbf{Uma limitação de plataforma confirmada:} no momento em que este
manual foi escrito, o módulo PowerShell do Hyper-V não expõe nenhum
cmdlet para registrar um certificado fornecido pelo usuário ---
\texttt{Set-VMFirmware -SecureBootTemplate} só aceita os templates
fixos da própria Microsoft. O VirtualBox é o backend a usar para uma
cadeia real de Secure Boot autoassinada hoje.

% =======================================================================
\chapter{Boot em Hardware Físico}
\label{ch:hardboot}
% =======================================================================

\section{A Diretiva HARDBOOT}

Por padrão, as imagens do cnimbus são construídas apenas para hardware
virtual: \texttt{virtio-net}, \texttt{virtio-blk}, e dispositivos
paravirtualizados similares que nenhuma máquina física realmente tem.
\dir{HARDBOOT} é uma diretiva opt-in, inerte por padrão, que troca por
drivers de hardware real em vez disso.

\begin{nimbusfile}
HARDBOOT eth      # ou: wifi, ou eth+wifi para ambos
\end{nimbusfile}

Isso também precisa ser passado para o \texttt{prepare} (muda qual
kernel é compilado):

\begin{shell}
cnimbus prepare --hardboot eth
\end{shell}

Um Nimbusfile cujo \dir{HARDBOOT} não corresponde às pieces contra as
quais está sendo construído é rejeitado em tempo de
\texttt{build-disk}, da mesma forma que um \dir{ARCH} incompatível já
é.

Cada valor isolado é exclusivo à sua própria família de drivers:
\dir{HARDBOOT eth} nunca constrói a pilha WiFi, e \dir{HARDBOOT wifi}
nunca constrói drivers de chipset Ethernet. \dir{HARDBOOT eth+wifi} é
o único valor que constrói ambos, para uma máquina física que tem
tanto uma NIC cabeada quanto um rádio WiFi.

\section{HARDBOOT eth}

Adiciona drivers reais de chipset Ethernet das famílias Intel
e1000/e1000e e Realtek R8169. \flag{FORMAT raw} é exigido para um boot
real via pendrive USB ou bare-metal (uma ISO simples não tem código de
boot MBR, então não pode ser escrita em um pendrive USB e inicializada
da forma que uma máquina real espera).

Como uma verificação prévia barata antes de gastar uma tentativa de
boot físico: a NIC emulada \textit{padrão} do QEMU e do VirtualBox é
um chip real da família e1000, então exatamente a mesma pilha de
drivers pode ser testada virtualmente primeiro
(\texttt{cnimbus run -{}-backend qemu -{}-format raw} ou
\texttt{-{}-backend vbox}). Isso não exercita o caminho real de
enumeração USB/MBR, mas pega de graça uma combinação obviamente
quebrada de driver/kconfig.

\section{HARDBOOT wifi}

Adiciona um conjunto curado de drivers de chipset WiFi (liderado pelo
Atheros AR9271, \texttt{ath9k\_htc}), os blobs de firmware
correspondentes (fixados por hash e embutidos no initramfs de boot), e
um \texttt{wpa\_supplicant} estaticamente ligado (o BusyBox não traz
nenhum). Use as diretivas \dir{WIFI} / \dir{WIFIPSK} / \dir{WIFICOUNTRY}
do Capítulo~\ref{ch:nimbusfile} para configurar a rede.

\textbf{Limitação honesta:} não há hardware de rádio WiFi, nem emulação
de chipset WiFi, disponível no QEMU ou VirtualBox --- diferente da
Ethernet cabeada, a prova de associação em hardware real deste perfil
genuinamente exige um adaptador WiFi físico. Tudo até esse ponto
(compilação do kernel, empacotamento de firmware, geração de
configuração, tratamento de credenciais) é verificado hoje por
inspeção direta de uma imagem real construída.

\section{HARDBOOT eth+wifi}

Constrói ambas as famílias de drivers no mesmo kernel e sobe ambas as
interfaces no boot (\texttt{eth0} via \dir{DHCP}/\dir{IP} como usual,
\texttt{wlan0} via \dir{WIFI}/\dir{WIFIPSK}/\dir{WIFICOUNTRY}) --- para
uma máquina física que genuinamente tem tanto uma porta cabeada quanto
um rádio WiFi. Exige as mesmas diretivas
\dir{WIFI}/\dir{WIFIPSK}/\dir{WIFICOUNTRY} que \dir{HARDBOOT wifi}
isoladamente exige; as mesmas duas lacunas de hardware real acima
(boot USB/BIOS, associação WiFi) se aplicam.

\begin{nimbusfile}
HARDBOOT eth+wifi
WIFI MyNetwork
WIFIPSK correcthorsebattery
WIFICOUNTRY BR
DHCP true
\end{nimbusfile}

% =======================================================================
\chapter{Exemplos Práticos}
\label{ch:examples}
% =======================================================================

Cada exemplo abaixo vive em \texttt{examples/} no repositório e é
totalmente autocontido: seu próprio Nimbusfile, seu próprio
\texttt{README.md} com comandos exatos de build/boot, e (quando um
binário complementar é necessário) seu próprio código-fonte pequeno em
Go. Todos eles assumem que o \texttt{cnimbus} está compilado e que o
\texttt{cnimbus prepare} já produziu \texttt{./pieces} para o \dir{ARCH}
de destino.

\section{env-config --- Passando Configuração em Tempo de Execução}

Mostra \dir{ENV}: valores de configuração embutidos em tempo de build,
lidos pelo entrypoint como variáveis de ambiente comuns. O exemplo mais
simples possível da filosofia ``declare, não script''.

\section{volume-persistent-disk --- Estado Que Sobrevive a um Reboot}

Mostra \dir{VOLUME}: montando um disco pré-formatado e pré-anexado, de
forma que um arquivo de log (ou qualquer outro estado) escrito durante
um boot ainda esteja lá após um ciclo completo de energia. Todo o resto
em uma imagem cnimbus é apenas-RAM por design --- esta é a única
válvula de escape deliberada.

\begin{nimbusfile}
VOLUME /dev/vda /mnt/data
ENTRYPOINT ["/bin/sh", "-c",
  "echo \"booted at $(date)\" >> /mnt/data/log.txt; sleep 3600"]
\end{nimbusfile}

\section{agent-http-kv --- Configuração ao Vivo Via HTTP Puro}

Mostra \dir{AGENT <url>}: uma VM em execução consultando um endpoint
HTTP do lado do host a cada poucos segundos e escrevendo a resposta em
\texttt{/var/run/cnimbus-kv.json}, permitindo que sua própria aplicação
capte mudanças de configuração sem um rebuild ou reboot. Funciona em
qualquer hypervisor com rede no guest.

\section{agent-vboxguest-kv --- Configuração ao Vivo Sem Rede no Guest}

A mesma ideia acima, mas pelo canal real de Guest Properties do
VirtualBox em vez de HTTP --- útil quando o guest deliberadamente não
tem acesso à rede nenhum, e a configuração ainda precisa chegar até
ele.

\section{multi-service --- Mais de Um Processo}

Mostra \dir{SERVICE}: executando um processo adicional, supervisionado
e reiniciado automaticamente, ao lado do \dir{ENTRYPOINT}, cada um
gerenciado independentemente.

\section{static-ip-firewall --- Rede Mais um Firewall}

Mostra \dir{IP} (um endereço estático, demonstrando que ele vence o
\dir{DHCP} quando ambos estão definidos) e \dir{FIREWALL}: aplicando
regras reais de \texttt{iptables} no boot.

\begin{nimbusfile}
IP 192.168.56.10 255.255.255.0 192.168.56.1
FIREWALL -P INPUT DROP
FIREWALL -A INPUT -p tcp --dport 8080 -j ACCEPT
\end{nimbusfile}

Tráfego de loopback e conexões estabelecidas/relacionadas são sempre
aceitas automaticamente antes de qualquer regra \dir{FIREWALL} rodar
--- você nunca precisa declarar nenhuma das duas você mesmo.

\section{format-raw-disk --- Imagens de Disco Estilo Nuvem}

Mostra \flag{FORMAT raw}: produzindo um disco raw particionado em GPT
em vez de uma ISO, o formato que plataformas de nuvem (Proxmox, e
provisionamento baseado em imagem de disco similares) esperam, e o
mesmo formato que \dir{HARDBOOT} exige para um boot físico.

\section{healthcheck-restart --- Detectando um Processo Travado}

Mostra \dir{HEALTHCHECK} e \dir{RESTART}: um servidor deliberadamente
instável que para de responder (sem travar) após 20 segundos, pego e
forçado a reiniciar por três verificações de saúde consecutivas
falhas.

\begin{nimbusfile}
HEALTHCHECK --interval=5 --retries=3 wget -q -O /dev/null http://127.0.0.1:8080/health
RESTART entrypoint always
\end{nimbusfile}

O console mostra a sequência exata de escalonamento: três verificações
falhas, depois \texttt{SIGTERM}, depois (se ainda estiver vivo)
\texttt{SIGKILL} (código de saída 137 = 128+9, confirmando que foi
mesmo um kill), depois um reinício limpo.

\section{stopgrace-graceful-shutdown --- Dando Tempo Para Trabalho em Andamento Terminar}

Mostra \dir{STOPGRACE}: um servidor que termina uma ``transação'' em
andamento de 8 segundos após receber \texttt{SIGTERM}, em vez de
morrer imediatamente --- o cenário para o qual o orçamento padrão de
cerca de 1 segundo de shutdown do init do BusyBox é muito curto.

\begin{nimbusfile}
STOPGRACE 15
\end{nimbusfile}

\section{hardboot-eth --- Ethernet Real, Boot USB Real}

Mostra \dir{HARDBOOT eth} de ponta a ponta: o passo
\flag{prepare --hardboot eth}, uma verificação prévia virtual com
e1000/e1000e sob o QEMU e o VirtualBox, e os passos exatos de
\texttt{dd}/Rufus para escrever a imagem raw resultante em um pendrive
USB para um boot físico real.

\section{hello-http-baremetal-proxmox --- Uma Imagem, Dois Destinos, IPv4 e IPv6}

Um único Nimbusfile que sobe sem modificação em hardware físico real
(via pendrive UEFI) e dentro do Proxmox (como o CD-ROM virtual de uma
VM), alcançável por qualquer lugar da internet tanto em IPv4 quanto em
IPv6:

\begin{nimbusfile}
KERNEL latest-stable
BUSYBOX latest
ARCH amd64
HOSTNAME cnimbus-hello-http-demo

WIFI cnimbus-placeholder-ssid
WIFIPSK cnimbus-placeholder-psk
WIFICOUNTRY US
HARDBOOT eth+wifi

VGA true
DHCP true

COPY ./helloserver /usr/bin/helloserver
ENTRYPOINT /usr/bin/helloserver
CMD :8080

FIREWALL -P INPUT DROP
FIREWALL -A INPUT -p tcp -s 0.0.0.0/0 --dport 8080 -j ACCEPT
FIREWALL6 -P INPUT DROP
FIREWALL6 -A INPUT -p tcp -s ::/0 --dport 8080 -j ACCEPT

FORMAT iso
\end{nimbusfile}

\dir{HARDBOOT eth+wifi} é aditivo sobre os drivers de VM que
\texttt{vm-amd64.fragment} já traz (e1000/virtio-net) --- é exatamente
isso que permite uma única imagem inicializar tanto em hardware real
quanto dentro do Proxmox sem mudança nenhuma: a NIC virtual do Proxmox
já é coberta pelos drivers de VM, independentemente do \dir{HARDBOOT};
a NIC (ou placa WiFi) de uma máquina real é coberta pelos drivers que o
\dir{HARDBOOT} adiciona por cima. \dir{WIFI}/\dir{WIFIPSK}/\dir{WIFICOUNTRY}
são obrigatórios sempre que \dir{HARDBOOT} inclui \texttt{wifi} --- os
valores acima são apenas placeholders (este exemplo só exercita
Ethernet cabeada nos dois destinos); substitua-os por credenciais reais
se o hardware físico de destino precisar associar via WiFi.

\dir{VGA true} ativa o banner de IP do console: assim que a rede sobe,
a tela (o monitor físico, ou o console noVNC/SPICE do Proxmox) mostra
algo como:

\begin{shell}
cnimbus: IPv4 address: 10.0.2.15 (eth0)
cnimbus: IPv6 address: fd00::... (eth0)
\end{shell}

\dir{FIREWALL}/\dir{FIREWALL6} (seção ``Referência Completa de
Diretivas'') são dois conjuntos de regras independentes --- declarar
um não afeta o outro. Ambos ficam abertos a \texttt{0.0.0.0/0}/
\texttt{::/0} na porta 8080 aqui, exatamente como pedido; loopback e
conexões já estabelecidas são sempre aceitas automaticamente antes de
qualquer regra rodar (mesmo comportamento do exemplo
\texttt{static-ip-firewall} anterior).

Verificado de ponta a ponta em hardware físico real e num Proxmox
real: o kernel sobe, o DHCP obtém um lease em \texttt{eth0}, o banner
de IP imprime tanto o endereço IPv4 quanto o IPv6, um \texttt{curl} de
outra máquina na mesma rede (inclusive \texttt{curl -6}) retorna a
resposta HTTP real, os dois conjuntos de \dir{FIREWALL}/\dir{FIREWALL6}
se aplicam sem nenhuma falha de regra, e os dois sinais de controle do
Proxmox (``Signal Shutdown'' via ACPI e Ctrl+Alt+Del pelo console)
desligam/reiniciam a VM de forma correta.

\textbf{Boot em hardware real (apenas UEFI --- esta ISO não é
isohybrid, veja ``Imagens ISO Não São Isohybrid''):}

\begin{shell}
# Linux/macOS
sudo dd if=cnimbus-hello-http-demo.iso of=/dev/sdX bs=4M \
  status=progress conv=fsync
# Windows: Rufus, modo "DD Image", apontado para o .iso
\end{shell}

Um pendrive gravado desta forma inicializa em qualquer hardware UEFI
com um disco USB anexado como um disco SATA/SCSI comum ---
\texttt{HARDBOOT eth}/\texttt{eth+wifi} traz o driver de
armazenamento USB necessário para o kernel redescobrir o próprio
pendrive depois que o firmware UEFI entrega o controle a ele. Uma
ferramenta de pendrive multiboot (Ventoy e equivalentes, que
inicializam encadeando o GRUB contra um arquivo \texttt{.iso} numa
partição FAT/exFAT em vez de gravar a ISO diretamente no disco) também
funciona sem configuração adicional: a imagem varre toda partição
FAT/exFAT do pendrive procurando qualquer arquivo \texttt{*.iso} e o
monta via loop.

Se o pendrive carrega mais de uma imagem cnimbus, o console mostra
qual delas realmente subiu --- lido de um \texttt{CNIMBUS.CFG} em
texto puro na raiz da própria ISO:

\begin{shell}
cnimbus: identified boot image:
HOSTNAME=cnimbus-hello-http-demo
ARCH=amd64
FORMAT=iso
CNIMBUS_VERSION=...
\end{shell}

\textbf{Boot no Proxmox:} envie a ISO para um storage que aceite
imagens ISO, crie uma VM sem disco algum (a imagem inteira roda em
RAM), anexe a ISO como unidade de CD/DVD, defina a ordem de boot para
inicializar a partir dela, e inicie a VM.

\section{Mensagens de Boot em Todos os Consoles}

Com \flag{--vga}/\dir{VGA true}, toda mensagem que a imagem imprime no
boot (o banner de IP, o estado dos serviços) aparece tanto no monitor
físico/console de vídeo quanto na porta serial, independente de qual
deles a linha de comando do kernel declara como \texttt{/dev/console}
principal. Isso importa porque o kernel sempre imprime seu próprio
\texttt{dmesg} em todos os consoles declarados, mas o userspace por
padrão escreveria apenas no último deles --- uma tela de monitor que
mostra o \texttt{dmesg} do kernel e nada da própria imagem não
significa que o boot travou, apenas que a saída está indo só para a
serial.

\section{Sinais de Controle do Proxmox: Signal Shutdown e Ctrl+Alt+Del}

Esta imagem responde aos dois sinais de controle que o Proxmox oferece
pela interface web:

\begin{itemize}
  \item \textbf{``Signal Shutdown''} envia um pedido de desligamento
    ACPI; a imagem completa o desligamento gracioso (para os serviços,
    desmonta o que precisa ser desmontado, e desliga a VM) igual a
    rodar \texttt{poweroff} localmente teria.
  \item \textbf{Ctrl+Alt+Del} pelo console reinicia a VM através do
    mesmo ciclo gracioso de \dir{RESTART}/\dir{STOPGRACE}.
\end{itemize}

Os dois foram verificados contra um Proxmox real e contra um Hyper-V
Generation~1 real.

\section{secureboot-uki --- Assinando e Registrando um Certificado Real}

Mostra \flag{--secureboot}/\flag{--uki} de ponta a ponta, incluindo os
comandos exatos e testados de \texttt{VBoxManage modifynvram} para
registrar um certificado real, autogerado, na variável \texttt{db}
UEFI de uma VM VirtualBox (usando o \texttt{cert-to-efi-sig-list} do
\texttt{efitools} para construir uma \texttt{EFI\_SIGNATURE\_LIST} real)
e inicializar com o Secure Boot habilitado --- o primeiro boot real
verificado por Secure Boot de um kernel assinado pelo cnimbus
documentado neste projeto, sem usar nenhuma chave de demonstração de
fornecedor.

% =======================================================================
\chapter{Solução de Problemas e Limitações Conhecidas}
\label{ch:troubleshooting}
% =======================================================================

\section{Nenhum Shell em Lugar Nenhum}

Não há shell instalado na imagem por design (além do que a própria
invocação em forma shell do \dir{ENTRYPOINT} usa internamente). Você
não pode fazer \texttt{ssh} ou login no console de uma VM cnimbus em
execução para explorar interativamente --- se você precisa de
visibilidade em tempo de execução, construa isso dentro do seu próprio
\dir{ENTRYPOINT}, ou use \dir{AGENT} para configuração ao vivo e um log
apoiado em VOLUME para saída persistida.

\section{Imagens ISO Não São Isohybrid}

Uma imagem \flag{FORMAT iso} simples tem entradas duplas de catálogo
de boot El~Torito (BIOS e UEFI em amd64) mas nenhum código de boot MBR
no byte~0. Ela vai inicializar corretamente na unidade óptica virtual
de qualquer hypervisor suportado, mas escrevê-la diretamente em um
pendrive USB e esperar que uma máquina real inicialize a partir dela
não vai funcionar. Use \flag{FORMAT raw} para boot USB/bare-metal em
vez disso --- esta é uma escolha de design explícita e documentada,
não um descuido.

\section{Aceleração WHPX no Windows}

\flag{--accel auto} / \flag{--accel whpx} foi observado falhando
completamente em pelo menos um host Windows real com
\texttt{QEMU: WHPX: Unexpected VP exit code 4}. \flag{--accel tcg}
(emulação por software) é um fallback confiável --- mais lento, mas
inicializa.

\section{VMware Player e vTPM}

O VMware Player (a versão gratuita) não consegue ligar uma VM com
\texttt{vtpm.present = "TRUE"} --- exige que a VM esteja criptografada
primeiro, um recurso exclusivo do Workstation Pro/vSphere. Use o
VirtualBox ou o Hyper-V se seu fluxo de trabalho precisa de um TPM
virtual.

\section{Arquivos Grandes em COPY/ADD e FORMAT iso}

\flag{FORMAT iso} tem um limite rígido em cima do \dir{TMPSIZE}: a
carga combinada de kernel + initramfs do estágio 1 é carregada via a
entrada de boot ``sem emulação'' do El~Torito, cujo campo de tamanho
tem 16 bits de unidades de 512 bytes (aproximadamente um orçamento
prático de 24~MiB para o initramfs comprimido). Um \dir{COPY}/\dir{ADD}
em um dos quatro diretórios executáveis em RAM grande o suficiente
para exceder isso faz o \texttt{build-disk} falhar com uma mensagem
nomeando o arquivo culpado. \flag{FORMAT raw} não tem esse limite ---
troque para ele para qualquer coisa grande.

\section{Códigos de Saída}

O \texttt{cnimbus} diferencia modos de falha via código de saída em vez
de sempre sair com 1:

\begin{itemize}
  \item \textbf{1} --- uma falha genérica/não classificada.
  \item \textbf{4} --- uma falha de integridade ou verificação de
    assinatura das pieces.
  \item \textbf{5} --- uma busca upstream (kernel.org, BusyBox, etc.)
    não pôde ser alcançada ou resolvida.
\end{itemize}

Execute \texttt{cnimbus help} para a lista autoritativa e atualizada.

% =======================================================================
\appendix
\chapter{Referência Rápida de Diretivas}
% =======================================================================

\begin{center}
\begin{longtable}{p{3.2cm}p{10.3cm}}
\toprule
\textbf{Categoria} & \textbf{Diretivas} \\
\midrule
\endhead
Alvo de boot & \dir{KERNEL}, \dir{BUSYBOX}, \dir{ARCH}, \dir{VGA}, \dir{HARDBOOT}, \dir{FORMAT} \\
Rede & \dir{HOSTNAME}, \dir{DHCP}, \dir{IP}, \dir{DNS}, \dir{NTP}, \dir{WIFI}, \dir{WIFIPSK}, \dir{WIFICOUNTRY} \\
Processo/SA & \dir{USER}, \dir{WORKDIR}, \dir{ENV}, \dir{LABEL}, \dir{EXPOSE}, \dir{ARG}, \dir{COPY}, \dir{ADD}, \dir{ENTRYPOINT}, \dir{CMD}, \dir{SERVICE} \\
Armazenamento & \dir{VOLUME}, \dir{TMPSIZE} \\
Config. ao vivo & \dir{AGENT} (\texttt{http}, \texttt{vboxguest}, \texttt{virtio-serial}, \texttt{aws-imds}, \texttt{ibm-imds}, \texttt{vmware}) \\
Firewall/ciclo de vida & \dir{FIREWALL}, \dir{FIREWALL\_ON\_ERROR}, \dir{HEALTHCHECK}, \dir{RESTART}, \dir{STOPGRACE} \\
Assinatura & \dir{PIECESKEY} \\
\bottomrule
\end{longtable}
\end{center}

% =======================================================================
\chapter{Métricas Medidas}
% =======================================================================

As métricas de requisitos não-funcionais a seguir foram medidas contra
um \texttt{prepare} + \texttt{build-disk} real, e (onde indicado) um
boot real em QEMU, na própria máquina de desenvolvimento deste projeto
--- veja \texttt{REQUIREMENTS.md} para a tabela completa e a
metodologia.

\begin{table}[h]
\centering
\begin{tabular}{lll}
\toprule
\textbf{Métrica} & \textbf{Medido} & \textbf{Orçamento} \\
\midrule
Tamanho da imagem (exemplo mínimo) & 19,11 MiB & 32 MiB \\
Tamanho do \texttt{vmlinuz} & 3,31 MiB & 12 MiB \\
RSS de pico do \texttt{build-disk} & 36,03 MiB & 512 MiB \\
Tempo de parede do \texttt{build-disk} (pieces em cache) & 0,53 s & 30 s \\
Boot até o entrypoint (emulação por software) & 6,79 s & --- \\
\bottomrule
\end{tabular}
\caption{Métricas medidas de NFR.}
\end{table}

\vfill
\begin{center}
\rule{0.5\textwidth}{0.4pt}\\[0.3cm]
\textit{Este manual resume o próprio README.md, BUILD.md, Tasks.md, e
REQUIREMENTS.md do projeto no momento em que foi escrito. Esses
arquivos, e os logs reais de boot referenciados ao longo do histórico
de commits deste projeto, permanecem a fonte de verdade autoritativa e
continuamente atualizada.}
\end{center}

\end{document}
